\documentclass[letterpaper]{article}

\usepackage[letterpaper,margin=1in]{geometry}

% Compile from the repository root so this path resolves to the pinned WSP
% submodule.
\input{wsp/testing/test-case-library.tex}

\begin{document}
\def\TCID{TC-0003}
\def\TCTitle{Socket layer over loopback}
\def\TCPurpose{Verify the non-blocking socket layer behind the Winsock
control with a real TCP connection and a UDP datagram on the loopback
address.}
\def\TCPriority{\TCPriorityHigh}

\def\TCRequirementRef{REQ-0001}
\def\TCDesignRef{Not applicable.}
\def\TCFeatureRef{\texttt{include/wfc/mswinsck/winsock\_io.hpp};
\texttt{src/winsock\_io.cpp}; \texttt{tests/winsock\_io\_tests.cpp}.}
\def\TCTestTechnique{Scenario testing of a TCP and a UDP exchange; each wait
polls up to five seconds.}

\def\TCPreconditions{Windows with WinSock 2; the loopback interface is
available; the test program \texttt{winsock\_io\_tests} is built (target
\texttt{wfc\_mswinsck\_io}).}
\def\TCEnvironment{Windows; registered only when \texttt{WIN32} is set. The
program binds to \texttt{127.0.0.1} with ports the system picks, so it does
not conflict with other programs. CTest runs it with a 60 second timeout.}
\def\TCAssumptions{The loopback interface delivers the data within five
seconds.}

\def\TCInputData{The text \texttt{hello} (TCP) and \texttt{datagram} (UDP).}
\def\TCInitialState{No socket is open.}

\def\TCProcedure{\begin{enumerate}
  \item Run \texttt{winsock\_io\_tests} (CTest \texttt{TC-0003-io}).
  \item TCP: open a server socket, bind it to port 0 and listen; the port it
    got is positive. Open a client, connect to it and wait until the
    connection completes. Accept the connection, check that the peer address
    is \texttt{127.0.0.1}, send \texttt{hello} from the client and receive it
    on the accepted socket. Close the client and check that the accepted
    socket reads end of stream. Close the remaining sockets.
  \item UDP: open and bind two sockets; send \texttt{datagram} from one to the
    port of the other and receive it, with the sender address and port.
  \item Check that the host name is not empty and that the text of socket
    error 10061 is not empty.
\end{enumerate}}
\def\TCExpectedResult{The program exits with status 0 and prints
\texttt{winsock io: tcp, udp and names are right}; no \texttt{FAIL} line is
written. The connection completes, the data arrives unchanged, the peer and
the sender are \texttt{127.0.0.1}, and the close is seen.}
\def\TCPostConditions{All sockets are closed; no files are written.}

\TCTemplate

\end{document}